\section{A metric keeps only the properties we use}\label{ch11:sec:metric}
Let $X$ be a set, whose members we call \emph{points}, even when they are numbers or vectors. A distance on $X$ should be a function $d$ from $X\times X$ to $\RR$ (Sections~\ref{ch03:sec:sets} and~\ref{ch03:sec:functions}).
\par\noindent\begin{minipage}{\linewidth}
\begin{definition}[Metric]\label{def:metric}
A \emph{metric} on a set $X$ assigns a real number $d(x,y)$ to every pair of points $x,y\in X$ and satisfies these requirements for all $x,y,z\in X$:
\begin{enumerate}
\item $d(x,y)\ge0$, and $d(x,y)=0$ holds exactly when $x=y$.
\item $d(x,y)=d(y,x)$.
\item $d(x,z)\le d(x,y)+d(y,z)$.
\end{enumerate}
A set together with a chosen metric is called a \emph{metric space} and is written $(X,d)$. The number $d(x,y)$ is called the \emph{distance} between $x$ and $y$.
\end{definition}
\end{minipage}\par

The first requirement says, in particular, that two different points are a positive distance apart. The second is called \emph{symmetry}, and the third is the triangle inequality of Section~\ref{ch11:sec:distance}, now demanded of $d$. To show that a new formula is a metric, we must \emph{recognize} it as one in the sense of Section~\ref{ch01:sec:definition}, by checking all three requirements; calling a formula a distance does not make it one.

\pauseq{On the real line, let $d(x,y)=(x-y)^2$. Is $d$ a metric?}{No. The first two requirements hold: a square is never negative, $(x-y)^2=0$ exactly when $x=y$, and $(y-x)^2=(x-y)^2$. The triangle inequality fails. Take $x=0$, $y=1$ and $z=2$. Then $d(x,z)=4$, while $d(x,y)+d(y,z)=1+1=2$, so the route through $y$ comes out \emph{shorter} than the direct route. One counterexample is enough.}

\subsection*{Examples}
The real line with $d(x,y)=|x-y|$ is a metric space. The absolute value is never negative, and $|x-y|=0$ exactly when $x=y$; symmetry and the triangle inequality were proved in Section~\ref{ch11:sec:distance}. When we treat the real line as a metric space without naming a metric, this is the metric we mean.

The set $\RR^n$ of vectors $x=(x_1,\ldots,x_n)$ from Section~\ref{ch10:sec:vectors} carries several useful metrics. The first is the \emph{Euclidean distance} $d_2(x,y)=\norm{x-y}_2$, the length of the difference vector; in the plane it is the straight-line distance, so from $(0,0)$ to $(3,4)$ it is $\sqrt{3^2+4^2}=5$. It is a metric: a length is never negative and is zero only for the zero vector (Section~\ref{ch10:sec:inner}), which $x-y$ is exactly when $x=y$; changing the sign of every coordinate of $x-y$ does not change the squares under the root, so $d_2(y,x)=d_2(x,y)$; and the triangle inequality $\norm{x-z}_2\le\norm{x-y}_2+\norm{y-z}_2$ is proved at the end of Section~\ref{ch10:sec:cs}.

Another useful choice is
\[
d_\infty(x,y)=\max_{1\le j\le n}|x_j-y_j|,
\]
the largest of the $n$ numbers $|x_1-y_1|,\ldots,|x_n-y_n|$; the subscript $\infty$ is just part of the traditional name. This metric records the worst error in any single coordinate. For $x=(1,5)$ and $y=(4,3)$, the coordinate differences have sizes $3$ and $2$, so $d_\infty(x,y)=3$, while $d_2(x,y)=\sqrt{13}\approx3.61$.

To prove the triangle inequality for $d_\infty$, let $x,y,z\in\RR^n$ and fix a coordinate $j$. The number $|x_j-y_j|$ is one of those whose largest is $d_\infty(x,y)$, so $|x_j-y_j|\le d_\infty(x,y)$; likewise $|y_j-z_j|\le d_\infty(y,z)$. With the triangle inequality for real numbers,
\[
|x_j-z_j|\le|x_j-y_j|+|y_j-z_j|\le d_\infty(x,y)+d_\infty(y,z).
\]
The right side does not depend on $j$, so it is at least each of the numbers $|x_j-z_j|$, and hence at least their largest: $d_\infty(x,z)\le d_\infty(x,y)+d_\infty(y,z)$. Exercise~\ref{ex:11.3} checks the other two requirements.

Which metric to use depends on the question. If an instrument must get \emph{every} coordinate right to within $0.01$, the condition is $d_\infty(x,y)\le0.01$; a bound on $d_2$ instead controls all the coordinate errors together, through the sum of their squares. The reading $x=(0.01,0.01)$ of the true point $y=(0,0)$ meets the first condition, since $d_\infty(x,y)=0.01$, but fails $d_2(x,y)\le0.01$, since $d_2(x,y)\approx0.0141$.

\subsection*{Norms, subsets, and spaces without addition}
Both metrics on $\RR^n$ measure the size of a difference vector: $d_2(x,y)=\norm{x-y}_2$, and $d_\infty(x,y)=\norm{x-y}_\infty$ with $\norm{x}_\infty=\max_j|x_j|$. Such measures of size are called norms. Precisely, a \emph{norm} on $\RR^n$ assigns to each vector $x$ a real number $\norm{x}$ such that, for all vectors $x,y$ and every real number $c$,
\begin{enumerate}
\item $\norm{x}\ge0$, with equality exactly when $x$ is the zero vector;
\item $\norm{cx}=|c|\,\norm{x}$;
\item $\norm{x+y}\le\norm{x}+\norm{y}$.
\end{enumerate}
The Euclidean length is a norm, by Sections~\ref{ch10:sec:inner} and~\ref{ch10:sec:cs}. So is $\norm{x}_\infty$: it is never negative and is $0$ exactly when every coordinate is $0$, multiplying every coordinate by $c$ multiplies the largest coordinate size by $|c|$, and the third property is proved one coordinate at a time, like the triangle inequality for $d_\infty$. Conversely, every norm gives a metric $d(x,y)=\norm{x-y}$: the first requirement comes from the first property of the norm, since $x-y$ is the zero vector exactly when $x=y$; symmetry from the second with $c=-1$, since $y-x=(-1)(x-y)$; and the triangle inequality from the third, since $x-z=(x-y)+(y-z)$.

A metric, however, needs only one number $d(x,y)$ for each pair of points, not a way to add points. So the rest of the chapter works in a general metric space, and its arguments apply also to sets in which addition makes no sense, or can lead outside the set. The most common such set is part of a larger metric space. If $(X,d)$ is a metric space and $C\subseteq X$, then $d$, used only for pairs of points of $C$ (the restriction of $d$ to $C\times C$, in the language of Section~\ref{ch03:sec:injsurj}), is a metric on $C$: the three requirements hold for all points of $X$, so in particular for all points of $C$. For example, the interval $[0,1]$ with distance $|x-y|$ is a metric space, even though a sum of its points can leave it, as $0.7+0.6=1.3$ does.
